\chapter*{Preface}
\addcontentsline{toc}{chapter}{Preface}

\section*{A profitable month and a bounced rent check}

Imagine a friend who runs a small software business from a rented studio. On the last afternoon of March, she calls you with a puzzle. Her team sold \$18,000 of subscriptions in March and earned \$4,000 of profit. Yet that same afternoon, her \$2,500 studio rent check bounced.

How can a business make \$4,000 of profit and still bounce a \$2,500 check? Trace every dollar in plain arithmetic. Suppose she started March with \$12,000 in her checking account. During March, customers signed \$18,000 of contracts while her costs totaled \$14,000 (\$11,500 for cloud servers and part-time help, plus \$2,500 for studio rent). Sales of \$18,000 minus costs of \$14,000 truly equals \$4,000 of profit. However, her contracts give business buyers thirty days to pay, so \$0 of March sales arrived in cash yet. Paying the \$11,500 server and wage bills drained her \$12,000 bank account down to just \$500 -- so when the landlord presented the \$2,500 rent check, it bounced.

Neither report is lying to her. Her business genuinely created \$4,000 of new value in March, even as her checking balance hit \$0. The confusion comes from mixing up two fundamentally different kinds of measurement.

\section*{The bathtub rule: stocks versus flows}

Picture an ordinary bathtub with a water tap above it and a drain at the bottom. You can ask two completely different questions about that tub.

First, freeze time at 8:00~PM and measure how much water sits in the tub right now. Suppose the tub holds $50$ liters at that exact instant. That snapshot at one moment in time is called a \textbf{stock} (an amount sitting in a container at a single timestamp). Your checking account balance at midnight tonight is a stock.

Second, unfreeze the clock and watch the tub from 8:00~PM to 8:10~PM. Suppose the tap pours in $6$ liters per minute while the drain lets out $4$ liters per minute. Over ten minutes, $60$ liters enter and $40$ liters leave, adding a net $20$ liters. A rate accumulated across a stretch of time is called a \textbf{flow} (how much enters or leaves over a period, such as dollars per month). Monthly sales, monthly salaries, and monthly rent are all flows.

The tub obeys one unbreakable rule. Water at 8:10~PM equals water at 8:00~PM, plus all water that poured in, minus all water that drained out:
\[
  50\text{ liters} + 60\text{ liters in} - 40\text{ liters out} = 70\text{ liters}.
\]

\begin{intuitionbox}
The bathtub rule clears up our friend's bounced rent check right away. Monthly profit measures a \emph{flow} of value earned across March, including bills customers have not paid yet. Bank cash is a \emph{stock} sitting in checking tonight. High monthly sales cannot pay Friday's rent unless actual cash has poured through the tap into the tub.
\end{intuitionbox}

If you have studied calculus, you already know a compact way to write this bathtub rule. Let $S(t)$ stand for the stock sitting in a bucket at instant $t$. Let $r(\tau)$ stand for the net flow rate (in dollars per day or per month) at moment $\tau$. Across the time interval from $t_0$ to $t_1$, the total flow accumulated is
\[
  F(t_0, t_1) = \int_{t_0}^{t_1} r(\tau)\,\dd\tau,
\]
and the ending stock equals the starting stock plus that accumulated flow:
\[
  S(t_1) = S(t_0) + F(t_0, t_1).
\]
If integrals look unfamiliar or rusty, do not worry. Business records slice time into discrete days, months, and quarters, turning that integral into plain addition and subtraction. Start Monday with \$100 in your wallet ($S(t_0) = 100$), earn \$50 at lunch, and spend \$30 at dinner. Your net flow is $F = 50 - 30 = 20$, so Monday ends at \$120 ($S(t_1) = 100 + 20 = 120$), as Figure~\ref{fig:preface-stock-flow} illustrates.

\begin{figure}[htbp]
\centering
\begin{tikzpicture}
  % Panel (a): Stock at instant t
  \begin{scope}[shift={(0,0)}]
    \node[anchor=south, font=\small\bfseries, text=sourceblue] at (3.20,4.85)
      {Panel (a): Stock Snapshot at Instant $t$};
    \draw[thick, gray!60] (0,0) -- (6.40,0);

    % Left stacked bar: What You Own (Assets)
    \filldraw[fill=sourceblue!18, draw=sourceblue, thick] (0.30,0) rectangle (2.75,1.80);
    \node[font=\footnotesize, align=center] at (1.525,0.90) {Cash \&\\Unpaid Bills};
    \filldraw[fill=sourceblue!32, draw=sourceblue, thick] (0.30,1.80) rectangle (2.75,3.60);
    \node[font=\footnotesize, align=center] at (1.525,2.70) {Laptops \&\\Servers};
    \node[font=\footnotesize\bfseries, text=sourceblue] at (1.525,-0.38) {Assets (Owned)};

    % Equals sign
    \node[font=\large\bfseries] at (3.20,1.80) {$=$};

    % Right stacked bar: Debts + Owners' Share
    \filldraw[fill=targetgreen!18, draw=targetgreen!80!black, thick] (3.65,0) rectangle (6.10,2.10);
    \node[font=\footnotesize, align=center] at (4.875,1.05) {Owners' Share\\(Equity)};
    \filldraw[fill=quotientorange!24, draw=quotientorange, thick] (3.65,2.10) rectangle (6.10,3.60);
    \node[font=\footnotesize, align=center] at (4.875,2.85) {Debts Owed\\(Liabilities)};
    \node[font=\footnotesize\bfseries, text=quotientorange!85!black] at (4.875,-0.38) {Who Claims It};

    \node[draw=sourceblue!70, fill=sourceblue!6, rounded corners=3pt, inner xsep=7pt, inner ysep=4pt, font=\footnotesize, align=center]
      at (3.20,4.22) {Frozen snapshot at moment $t$ (\$)};
  \end{scope}

  % Divider
  \draw[gray!40, dashed, thick] (7.15,-0.65) -- (7.15,5.25);

  % Panel (b): Flow across [t_0, t_1]
  \begin{scope}[shift={(8.05,0)}]
    \node[anchor=south, font=\small\bfseries, text=targetgreen!80!black] at (3.10,4.85)
      {Panel (b): Flow Accumulated Over $[t_0, t_1]$};
    \draw[->, thick, >=Latex] (-0.2,0) -- (6.0,0) node[right, font=\footnotesize] {Time};
    \draw[->, thick, >=Latex] (0,-0.2) -- (0,3.95) node[above, font=\footnotesize] {Rate (\$/mo)};

    % Shaded integral between t_0 = 0.95 and t_1 = 4.95
    \fill[targetgreen!20] (0.95,0) -- plot[domain=0.95:4.95, smooth] (\x, {1.40 + 0.48*sin(68*\x) + 0.18*\x}) -- (4.95,0) -- cycle;
    \draw[targetgreen!80!black, very thick, domain=0.45:5.35, smooth] plot (\x, {1.40 + 0.48*sin(68*\x) + 0.18*\x});

    \draw[dashed, gray!80] (0.95,0) -- (0.95,2.00);
    \draw[dashed, gray!80] (4.95,0) -- (4.95,2.10);
    \node[font=\footnotesize] at (0.95,-0.38) {Start $t_0$};
    \node[font=\footnotesize] at (4.95,-0.38) {End $t_1$};

    \node[font=\footnotesize, align=center, text=targetgreen!75!black] at (2.95,0.98)
      {Total Net Flow\\Added Over $[t_0, t_1]$};

    \node[draw=targetgreen!75, fill=targetgreen!8, rounded corners=3pt, inner xsep=7pt, inner ysep=4pt, font=\footnotesize, align=center]
      at (3.30,3.45) {$\text{Stock}(t_1) = \text{Stock}(t_0) + \text{Net Flow}$};
  \end{scope}
\end{tikzpicture}
\caption{Stocks versus flows. Left: At a single instant $t$, a Balance Sheet takes a snapshot of what the business owns and who has claims on it. Right: Across a time window from $t_0$ to $t_1$, flows of sales, costs, and payments move the business from its starting snapshot to its ending snapshot.}
\label{fig:preface-stock-flow}
\end{figure}

\begin{recoverybox}
\textbf{Pause and think:} Test the distinction on three everyday numbers before reading on. Consider (1) an \$84,000 yearly salary, (2) \$1,900 owed on a credit card tonight, and (3) \$650 spent on groceries in April. Which are stocks and which are flows? The \$1,900 card balance is a stock measured at one instant. The \$84,000 per year and \$650 per month are flows measured over a period.
\end{recoverybox}

\section*{Two simple ideas behind the whole subject}

From the outside, accounting and finance can look like a maze of tax forms and three-letter acronyms. Strip away the industry slang and only two core ideas run the entire subject.

First, \textbf{accounting tracks every dollar's origin and destination so nothing vanishes silently}. Whenever a business raises money, buys laptops, ships customer work, or pays salaries, \textbf{double-entry bookkeeping} (recording both where value came from and where it went) keeps the system in exact balance. Money cannot appear out of nowhere, and costs cannot disappear between the cracks.

Second, \textbf{finance compares money across time and risk}. Accounting tells you clearly what happened up to today. Finance helps you make forward-looking choices. Should you spend \$10,000 today on a machine that saves \$4,000 a year for three years? Why is \$100 in your hand today worth more than a promise of \$100 next year? Cash today can earn interest right away, whereas a future promise carries waiting time and uncertainty. Finance gives us a fair way to translate future dollars into today's dollars before comparing them.

\begin{focusbox}
Two rules guide every chapter ahead. First, \emph{double-entry accounting is a conservation law}: every transaction moves value from a clear source to a clear destination. Second, \emph{finance prices time and risk}: dollars at different dates or risk levels must be translated to today's value before you add or compare them.
\end{focusbox}

\section*{Where we are going}

This book assumes zero prior background in accounting, economics, or business. If you can read basic algebra and follow a table of small numbers, you have everything you need. We build upward in five connected parts:

\begin{itemize}
  \item \textbf{Part~I: Double-Entry Accounting and the Three Financial Statements (Chapters~1--2).} We follow two software engineers, Maya and Ben, as they launch a small hosting company from scratch. Starting from a plain notebook, we discover why double-entry bookkeeping works and build the three universal business reports: the Income Statement, the Balance Sheet, and the Cash Flow Statement.
  \item \textbf{Part~II: Evaluating Businesses, Public Filings, S-1s, and Multiples (Chapters~3--4).} We learn how interest compounds, how to translate future cash flows into present value, and how investors judge whether a stock or private business is cheap or expensive. Then we read real public company reports line by line.
  \item \textbf{Part~III: Startups, Fundraising, Cap Tables, and Founder Contracts (Chapters~5--6).} We trace how startups raise money, split ownership shares on a \textbf{cap table} (the table listing who owns what slice of a company), and divide the cash when a company is acquired.
  \item \textbf{Part~IV: Cross-Border Mobile Money, FX, Stablecoins, and Bank Reporting (Chapters~7--10).} We step inside an operating payment company that moves money across mobile wallets, foreign currencies, and digital stablecoins. We see how safeguarding customer funds and running daily three-way checks catch missing money before small bugs become disasters.
  \item \textbf{Part~V: Personal Taxation, Wrappers, and Individual Investing (Chapter~11).} We turn the same compounding and tax math onto your own paycheck, tax-sheltered savings accounts, government bonds, and long-term index funds.
\end{itemize}

\section*{Check the units before trusting any ratio}

Just like physics, finance obeys unit checks. Before you trust any financial ratio, write down the units on top and bottom:

\begin{itemize}
  \item \textbf{Stocks} are measured in \textbf{dollars} at a moment in time (for example, \$12,000 in checking on June~30, or \$500 in a wallet).
  \item \textbf{Flows} are measured in \textbf{dollars per unit of time} (for example, \$3,000 per month of spending, or \$36,000 per year).
  \item \textbf{Stock divided by flow gives time.} Divide \$12,000 of bank cash by \$3,000 per month of spending:
  \[
    \frac{\$12{,}000}{\$3{,}000\text{ / month}} = 4\text{ months}.
  \]
  The dollars cancel out, telling you the company has $4$~months of cash runway left. Valuation multiples in Part~II work the exact same way. Dividing a \$10~million business price (dollars) by \$1~million of yearly cash profit (dollars per year) gives $10$~years -- how long steady profits would take to earn back the purchase price.
  \item \textbf{Flow divided by stock gives a rate per time}, such as earning \$50 per year on a \$1,000 deposit ($5\%$ per year). \textbf{Flow divided by flow gives a pure fraction}, such as keeping \$20 of profit out of \$100 of monthly sales ($20\%$).
\end{itemize}

\begin{trapbox}
Unit and claim checks catch broken formulas right away. Suppose someone divides the total value of a business (which belongs to both bank lenders and share owners) by profit left over for share owners alone. Top and bottom belong to different groups of people, so the ratio is distorted before you even plug in numbers.
\end{trapbox}

\begin{happeningbox}
Throughout the book, short shaded boxes pause the main story to help ideas click. \emph{Intuition} boxes give plain-English mental models. \emph{Pause} boxes invite you to test a tiny number in your head before reading further. \emph{What is really happening} boxes strip away formal labels to show the underlying cash mechanics, while \emph{Common trap} and \emph{History} boxes highlight classic pitfalls and where key conventions came from.
\end{happeningbox}

\begin{intuitionbox}
Keep a scratchpad or a Python terminal open beside you as you read. Check every small table sum yourself, trace where each dollar goes, and work the short exercises at the end of each chapter. Once you see the simple conservation rules underneath the reports, you will never have to memorize an accounting rule again.
\end{intuitionbox}
